\begin{pastexam}[label={px:ncku-114-prog-03}]{成大資訊聯招 114 學年度 第 3 題}
（單選）

{[3\%]} The following C code is provided. Analyze the functionality of the code and answer the question.
\begin{minted}[fontsize=\small]{c}
#include <stdio.h>
int operation1(int parent[], int x) {
    if (parent[x] != -1) {
        int rootX = operation1(parent, parent[x]);
        parent[x] = rootX;
        return rootX;
    }
    return x;
}
void operation2(int parent[], int x, int y) {
    int rootX = operation1(parent, x);
    int rootY = operation1(parent, y);
    if (rootX != rootY) {
        parent[rootY] = rootX;
    }
}
int main() {
    int parent[6] = {-1, -1, -1, -1, -1, -1};
    parent[2] = 1;
    parent[3] = 2;
    parent[4] = 3;

    operation1(parent, 4);
    operation2(parent, 1, 3);
    operation2(parent, 5, 4);

    for(int i=0; i<=5;i++){
        printf("%d ", parent[i]);
    }
    return 0;
}
\end{minted}
Which of the following statements is correct?
\begin{enumerate}[label=(\Alph*)]
\item In the function main(), operation1(parent, 4) updates all nodes in the chain starting from 4, and operation2(parent, 1, 3) merges two different groups.
\item The output of the program is ``-1 -1 1 1 1 1''.
\item The output of the program is ``-1 5 1 2 3 -1''.
\item The function operation1() compresses paths, and the function operation2() merges groups if they have different representatives.
\item The function operation1() creates a loop, and the function operation2() merges all elements into one group.
\end{enumerate}
\end{pastexam}

\begin{pastsoln}
答案：(D)。

程式用 $-1$ 表示根（資結書用 $p[x] = x$，第~13.2.4~節的作答提醒提過用負數標記根的寫法）；\texttt{operation1}
就是資結書虛擬碼~13.5 的遞迴路徑壓縮，\texttt{operation2} 是資結書虛擬碼~13.2 不加任何合併規則的
\proc{Union}：$y$ 的根接到 $x$ 的根底下。 下面把 \texttt{operation1(parent, x)} 簡寫成 op1$(x)$，
\texttt{operation2(parent, x, y)} 簡寫成 op2$(x, y)$。 元素是 $0$ 到 $5$，每一步之後的 \texttt{parent[0..5]}：
\begin{center}
\small
\begin{tabular}{ll}
\toprule
步驟 & \texttt{parent[0..5]} \\
\midrule
設定 \texttt{parent[2..4]} 之後 & $-1,\ -1,\ 1,\ 2,\ 3,\ -1$（鏈 $4 \to 3 \to 2 \to 1$） \\
op1$(4)$ & $-1,\ -1,\ 1,\ \mathbf{1},\ \mathbf{1},\ -1$ \\
op2$(1, 3)$ & $-1,\ -1,\ 1,\ 1,\ 1,\ -1$（兩個根都是 $1$，不合併） \\
op2$(5, 4)$ & $-1,\ \mathbf{5},\ 1,\ 1,\ 1,\ -1$ \\
\bottomrule
\end{tabular}
\end{center}
op1$(4)$ 一路遞迴到根 $1$，回來時把 $4$、$3$、$2$ 的父節點都設成 $1$（$2$ 本來就是）。 op2$(1, 3)$ 找到的兩個根
都是 $1$，什麼都不做。 op2$(5, 4)$ 的兩個根是 $5$ 和 $1$，執行 \texttt{parent[1] = 5}。 程式印出
\texttt{-1 5 1 1 1 -1}（每個數後面一個空格）。
\begin{itemize}
\item (A) 錯。 前半句對（$4$ 到根的路徑上每個節點都改接到 $1$），但 op2$(1, 3)$ 的 $1$、$3$ 已經在同一組，沒有合併。
\item (B) 錯。 實際輸出是 \texttt{-1 5 1 1 1 -1}：$1$ 被接到 $5$ 底下，$5$ 仍然是根。
\item (C) 錯。 這是\textbf{不做}路徑壓縮時的輸出：\texttt{parent[3]}、\texttt{parent[4]} 會停在 $2$、$3$。
\texttt{operation1} 有壓縮，所以它們都變成 $1$。
\item (D) 對。 \texttt{operation1} 在遞迴回來時把路上每個節點直接接到根，就是路徑壓縮；\texttt{operation2} 先找
兩個代表，不同才合併。
\item (E) 錯。 \texttt{operation1} 只把節點接到根，不會形成迴圈；最後 $0$ 自己一組，$\{1, 2, 3, 4, 5\}$ 一組，並沒有
全部合成一組。
\end{itemize}
\end{pastsoln}
